%!TEX root = main.tex

\section{Introduction}

Unsupervised visual representation learning, or self-supervised learning, aims at obtaining features without using manual annotations and is rapidly closing the performance gap with supervised pretraining in computer vision~\cite{chen2020simple,he2019momentum,misra2019self}.
Many recent state-of-the-art methods build upon the instance discrimination task that considers each image of the dataset (or ``instance'') and its transformations as a separate class~\cite{dosovitskiy2016discriminative}. 
This task yields representations that are able to discriminate between different images, while achieving some invariance to image transformations. 
Recent self-supervised methods that use instance discrimination rely on a combination of two elements: (i) a contrastive loss~\cite{hadsell2006dimensionality} and (ii) a set of image transformations. 
The contrastive loss removes the notion of instance classes by directly comparing image features while the image transformations define the invariances encoded in the features.
Both elements are essential to the quality of the resulting networks~\cite{misra2019self,chen2020simple} and our work improves upon both the objective function and the transformations.

The contrastive loss explicitly compares pairs of image representations to push away representations from different images while pulling together those from transformations, or views, of the same image.
Since computing all the pairwise comparisons on a large dataset is not practical, most implementations approximate the loss by reducing the number of comparisons to random subsets of images during training~\cite{chen2020simple,he2019momentum,wu2018unsupervised}.
An alternative to approximate the loss is to approximate the task---that is to relax the instance discrimination problem.
For example, clustering-based methods discriminate between groups of images with similar features instead of individual images~\cite{caron2018deep}. 
The objective in clustering is tractable, but it does not scale well with the dataset as it requires a pass over the entire dataset to form image ``codes'' (\ie, cluster assignments) that are used as targets during training.
In this work, we use a different paradigm and propose to compute the codes online while enforcing consistency between codes obtained from views of the same image.
Comparing cluster assignments allows to contrast different image views while not relying on explicit pairwise feature comparisons.
Specifically, we propose a simple ``swapped'' prediction problem where we predict the code of a view from the representation of another view.
We learn features by \textbf{Sw}apping \textbf{A}ssignments between multiple \textbf{V}iews of the same image (\textbf{\OURS}).
The features and the codes are learned online, allowing our method to scale to potentially unlimited amounts of data.
In addition, \OURS works with small and large batch sizes and does not need a large memory bank~\cite{wu2018unsupervised} or a momentum encoder~\cite{he2019momentum}.

Besides our online clustering-based method, we also propose an improvement to the image transformations.
Most contrastive methods compare one pair of transformations per image, even though there is evidence that comparing more views during training improves the resulting model~\cite{misra2019self}. 
In this work, we propose \texttt{multi-crop} that uses smaller-sized images to increase the number of views while not increasing the memory or computational requirements during training. 
We also observe that mapping small parts of a scene to more global views significantly boosts the performance.
Directly working with downsized images introduces a bias in the features~\cite{touvron2019fixing}, which can be avoided by using a mix of different sizes. 
Our strategy is simple, yet effective, and can be applied to many self-supervised methods with consistent gain in performance.

We validate our contributions by evaluating our method on several standard self-supervised benchmarks.
In particular, on the ImageNet linear evaluation protocol, we reach $75.3\%$ top-$1$ accuracy with a standard ResNet-50, and $78.5\%$ with a wider model.
We also show that our \texttt{multi-crop} strategy is general, and improves the performance of different self-supervised methods, namely SimCLR~\cite{chen2020simple}, DeepCluster~\cite{caron2018deep}, and SeLa~\cite{asano2019self}, between $2\%$ and $4\%$ top-1 accuracy on ImageNet.
%Finally, we show the potential of our approach at a larger scale, by pre-training models on $1$ billion uncurated images from Instagram. 
%This pre-training improves the performance of supervised models by $0.9\%$ on ImageNet.
Overall, we make the following contributions:
\begin{itemize}[leftmargin=*]
\itemsep0.05em 

\item
We propose a scalable online clustering loss that improves performance by $+2\%$ on ImageNet and works in both large and small batch settings without a large memory bank or a momentum encoder.
\item
We introduce the \texttt{multi-crop} strategy to increase the number of views of an image with no computational or memory overhead.
We observe a consistent improvement of between $2\%$ and $4\%$ on \ImNet with this strategy on several self-supervised methods.
\item
Combining both technical contributions into a single model, we improve the performance of self-supervised by $+4.2\%$ on ImageNet with a standard ResNet and outperforms supervised ImageNet pretraining on multiple downstream tasks.
This is the first method to do so without finetuning the features, \ie, only with a linear classifier on top of frozen features.
%Besides, we show that pre-training a ResNet-$50$ on $1$ billion uncurated images from Instagram allows to reach a top-$1$ accuracy of $77.4\%$, which is an improvement of $+0.9\%$ over the same model trained from scratch.
\end{itemize}
